\section{Auslander--Buchsbaum}
\label{mk-chap:Albanese_AuslanderBuchsbaum}


\section{Depth of a module over a local ring}
\label{mk-sec:ab_depth}

\begin{definition}[Depth of a module]\label{mk-def:depth}
  \lean{RingTheory.Module.depth}
  \leanok
  \dcref{00LF,custom}

  \textit{Source: \cite{stacks-project}, tag 00LF (definition-depth).}
  Let \(R\) be a commutative ring, let \(I \subseteq R\) be an ideal, and let \(M\) be a finitely
  generated \(R\)-module. The \emph{\(I\)-depth} of \(M\), written \(\text{depth}_I(M)\), is the
  supremum in \(\{0, 1, 2, \dots, \infty\}\) of the lengths of \(M\)-regular sequences contained
  in \(I\), provided \(IM \neq M\); if \(IM = M\) we set \(\text{depth}_I(M) = \infty\). When
  \((R, \mathfrak{m})\) is a Noetherian local ring we abbreviate
  \(\text{depth}(M) := \text{depth}_{\mathfrak{m}}(M)\) and call this simply the
  \emph{depth} of \(M\). The depth of the ring \(R\) itself is \(\text{depth}(R)\), the depth
  of \(R\) as a module over itself.
\end{definition}

Recall that a sequence \(f_1, \dots, f_r \in R\) is called \emph{\(M\)-regular} if
\(M/(f_1, \dots, f_r)M \neq 0\) and for each \(i\), the element \(f_i\) is a non-zero-divisor on
\(M/(f_1, \dots, f_{i-1})M\). The empty sequence is regular on \(M\) iff \(M \neq 0\) (it is
\emph{not} regular on the zero module, but the convention \(\text{depth}(0) = \infty\)
above is the usual one).

The two practical computations of depth are the regular-sequence definition (used to
establish a \emph{lower} bound) and the \(\text{Ext}\) characterisation (used to establish an
\emph{upper} bound and to prove the depth-on-a-short-exact-sequence inequality).

\begin{lemma}[Depth via \(\text{Ext}\)]\label{mk-lem:depth_via_ext}
  \lean{RingTheory.Module.depth_eq_smallest_ext_index}
  \leanok
  \dcref{00LP,custom}

  \uses{mk-def:depth, mk-lem:ext_smul_eq_zero_of_mem_annihilator}

  \textit{Source: \cite{stacks-project}, tag 00LP (lemma-depth-ext).}
  Let \((R, \mathfrak{m})\) be a Noetherian local ring with residue field
  \(\kappa = R/\mathfrak{m}\) and let \(M\) be a nonzero finitely generated \(R\)-module. Then
  \[
    \text{depth}(M) \;=\; \inf\bigl\{\, i \in \mathbb{N} \;\bigm|\; \text{Ext}^i_R(\kappa, M) \neq 0 \,\bigr\}.
  \]
\end{lemma}

\begin{proof}
  \leanok


Write \(\delta := \text{depth}(M)\) and \(i := \inf\{ j : \text{Ext}^j_R(\kappa, M) \neq 0 \}\).
  The base case \(\delta = 0\) is exactly \(\text{Hom}_R(\kappa, M) \neq 0\) (some element of
  \(M\) is annihilated by \(\mathfrak{m}\), i.e.\ \(\mathfrak{m} \in \text{Ass}(M)\), i.e.\
  \(\mathfrak{m}\) contains a zero-divisor on \(M\)), which is \(i = 0\). For the inductive step
  pick \(x \in \mathfrak{m}\) a non-zero-divisor on \(M\) with
  \(\text{depth}(M/xM) = \delta - 1\). The short exact sequence
  \(0 \to M \xrightarrow{x} M \to M/xM \to 0\) induces a long exact sequence of \(\text{Ext}^*_R(\kappa,
  -)\). Multiplication by \(x\) annihilates each \(\text{Ext}^j_R(\kappa, M)\) (because
  \(x \in \mathfrak{m}\) and \(\kappa\) is annihilated by \(\mathfrak{m}\)), so the long exact
  sequence breaks into short pieces giving
  \(\text{Ext}^{j-1}_R(\kappa, M/xM) \cong \text{Ext}^j_R(\kappa, M)\) at the relevant
  indices, and the smallest non-vanishing index drops by \(1\). By induction
  \(i(M/xM) = i - 1 = \delta - 1\), so \(i = \delta\).
\end{proof}

\section{Projective dimension}
\label{mk-sec:ab_pd}

\begin{definition}[Projective dimension]\label{mk-def:projective_dimension}
  \lean{Module.projectiveDimension}
  \leanok
  \dcref{custom}

  Let \(R\) be a commutative ring and let \(M\) be an \(R\)-module. The \emph{projective
  dimension} of \(M\), written \(\text{pd}_R(M)\), is the infimum in
  \(\{0, 1, 2, \dots, \infty\}\) of the lengths \(n\) of projective resolutions
  \[
    0 \longrightarrow P_n \longrightarrow P_{n-1} \longrightarrow \cdots
    \longrightarrow P_1 \longrightarrow P_0 \longrightarrow M \longrightarrow 0
  \]
  with each \(P_i\) a projective \(R\)-module. We say \(M\) has \emph{finite projective dimension}
  if \(\text{pd}_R(M) < \infty\), and \emph{infinite projective dimension} otherwise.
  This chapter adopts the convention \(\text{pd}_R(0)=0\), represented by the empty
  resolution \(0 \to 0\); some references instead use \(-\infty\). The
  Auslander--Buchsbaum formula in \ref{mk-thm:auslander_buchsbaum} concerns only nonzero \(M\).
  In the local Noetherian setting a finitely generated \(R\)-module of finite projective
  dimension admits a \emph{minimal finite free resolution}: a resolution with each \(P_i\)
  free of finite rank and each transition map \(P_{i+1} \to P_i\) landing in \(\mathfrak{m} P_i\).
\end{definition}

The two basic facts about \(\text{pd}_R(M)\) needed downstream are:
\begin{itemize}
  \item \(\text{pd}_R(M) = 0\) iff \(M\) is projective (in the local Noetherian setting,
    iff \(M\) is free);
  \item if \(0 \to K \to P \to M \to 0\) is a short exact sequence with \(P\) projective, then
    \(\text{pd}_R(M) = \text{pd}_R(K) + 1\) provided \(K \neq 0\).
\end{itemize}
Both statements follow by splitting or splicing finite projective resolutions.

\section{Depth on a short exact sequence}
\label{mk-sec:ab_depth_ses}

\begin{lemma}[Depth on a short exact sequence]\label{mk-lem:depth_short_exact_sequence}
  \lean{RingTheory.Module.depth_of_short_exact}
  \leanok
  \dcref{00LE,custom}

  \uses{mk-lem:natCast_add_one_le_of_le_sub_one}

  \textit{Source: \cite{stacks-project}, tag 00LE (lemma-depth-in-ses).}
  Let \((R, \mathfrak{m})\) be a Noetherian local ring and let
  \(0 \to N' \to N \to N'' \to 0\) be a short exact sequence of nonzero finitely generated
  \(R\)-modules. Then
  \begin{enumerate}
    \item \(\text{depth}(N) \;\geq\; \min\bigl(\text{depth}(N'),\; \text{depth}(N'')\bigr)\);
    \item \(\text{depth}(N'') \;\geq\; \min\bigl(\text{depth}(N),\; \text{depth}(N') - 1\bigr)\);
    \item \(\text{depth}(N') \;\geq\; \min\bigl(\text{depth}(N),\; \text{depth}(N'') + 1\bigr)\).
  \end{enumerate}
\end{lemma}

\begin{proof}
  \leanok


\uses{mk-lem:depth_via_ext, mk-lem:ext_vanish_of_natCast_lt_depth, mk-lem:natCast_add_one_le_of_le_sub_one}
  By \ref{mk-lem:depth_via_ext}, \(\text{depth}(N')\) (resp.\ \(\text{depth}(N)\),
  \(\text{depth}(N'')\)) is the smallest index at which \(\text{Ext}^*_R(\kappa, -)\) is
  nonzero on \(N'\) (resp.\ \(N\), \(N''\)). Apply the long exact \(\text{Ext}^*_R(\kappa, -)\)
  sequence to \(0 \to N' \to N \to N'' \to 0\):
  \[
    \cdots \to \text{Ext}^{i-1}_R(\kappa, N'')
    \to \text{Ext}^i_R(\kappa, N')
    \to \text{Ext}^i_R(\kappa, N)
    \to \text{Ext}^i_R(\kappa, N'')
    \to \text{Ext}^{i+1}_R(\kappa, N') \to \cdots
  \]
  Each of the three inequalities is then a direct read-off of the fact that if two
  of the three modules in a length-3 chunk vanish then the middle vanishes (for (1)) /
  the boundary terms vanish (for (2),(3)). Compare with the diagonal vanishing argument in
  the proof of \ref{mk-lem:depth_via_ext}.
\end{proof}

\section{The Auslander--Buchsbaum formula}
\label{mk-sec:ab_main}

\begin{theorem}[Auslander--Buchsbaum formula]\label{mk-thm:auslander_buchsbaum}
  \lean{RingTheory.auslander_buchsbaum_formula}
  \leanok
  \dcref{090V,custom}


  \textit{Source: \cite{stacks-project}, tag 090V (proposition-Auslander-Buchsbaum); cf.\
  Matsumura, \emph{Commutative Ring Theory}, Theorem 19.1; Auslander--Buchsbaum,
  \emph{Homological dimension in local rings}, 1957.}
  Let \((R, \mathfrak{m})\) be a Noetherian local ring and let \(M\) be a nonzero finitely
  generated \(R\)-module of finite projective dimension. Then
  \[
    \text{pd}_R(M) \;+\; \text{depth}(M) \;=\; \text{depth}(R).
  \]
\end{theorem}

\begin{proof}
  \leanok


\uses{mk-lem:depth_short_exact_sequence, mk-lem:auslander_buchsbaum_formula_succ_pd,
    mk-lem:depth_eq_of_linearEquiv, mk-lem:depth_pi_const_eq_depth_of_nonempty}
  We induct on \(\text{depth}(M)\).

  \emph{Base case \(\text{depth}(M) = 0\).} Let \(e := \text{pd}_R(M)\) and pick a minimal
  finite free resolution
  \[
    0 \longrightarrow R^{n_e} \longrightarrow R^{n_{e-1}} \longrightarrow
    \cdots \longrightarrow R^{n_0} \longrightarrow M \longrightarrow 0
  \]
  with all transition maps having matrix coefficients in \(\mathfrak{m}\) (this is the
  ``minimal-resolution-after-trimming-trivial-summands'' normalisation, cf.\ Stacks
  \texttt{lemma-add-trivial-complex}). On the one hand, the ``what is exact'' criterion
  (Stacks tag 00MF, \texttt{proposition-what-exact}, expressing exactness of a sequence of
  finite free modules in terms of depths of ideals of \(r\)-minors) gives
  \(\text{depth}(R) \geq e\). On the other hand, splitting the resolution into short exact
  sequences
  \[
    0 \to R^{n_e} \to R^{n_{e-1}} \to K_{e-2} \to 0, \quad
    0 \to K_{e-2} \to R^{n_{e-2}} \to K_{e-3} \to 0, \quad \dots, \quad
    0 \to K_0 \to R^{n_0} \to M \to 0,
  \]
  and applying \ref{mk-lem:depth_short_exact_sequence} iteratively,
  yields the chain of inequalities
  \(\text{depth}(K_{e-2}) \geq \text{depth}(R) - 1\),
  \(\text{depth}(K_{e-3}) \geq \text{depth}(R) - 2\), \dots,
  \(\text{depth}(K_0) \geq \text{depth}(R) - (e-1)\), and finally
  \(\text{depth}(M) \geq \text{depth}(R) - e\). Since \(\text{depth}(M) = 0\) we deduce
  \(\text{depth}(R) \leq e\). Combined with the previous inequality, $\text{depth}(R) = e =
  \text{pd}_R(M) + 0 = \text{pd}_R(M) + \text{depth}(M)$.

  \emph{Inductive step \(\text{depth}(M) > 0\).} We claim there exists \(x \in \mathfrak{m}\)
  which is a non-zero-divisor on both \(R\) and \(M\). Indeed: either \(\text{pd}_R(M) > 0\), in
  which case \(\text{depth}(R) > 0\) by the ``what is exact'' criterion applied to the start
  of a minimal resolution, or \(\text{pd}_R(M) = 0\) in which case \(M\) is finite free and
  \(\text{depth}(R) = \text{depth}(M) > 0\). Either way both \(\text{depth}(R) > 0\) and
  \(\text{depth}(M) > 0\), hence \(\text{depth}(R \oplus M) > 0\) by
  \ref{mk-lem:depth_short_exact_sequence}, so a common non-zero-divisor \(x \in
  \mathfrak{m}\) exists. Take a minimal finite free resolution of \(M\) as above; the snake
  lemma applied to multiplication by \(x\) gives a minimal finite free resolution
  \[
    0 \to (R/xR)^{n_e} \to (R/xR)^{n_{e-1}} \to \cdots \to (R/xR)^{n_0} \to M/xM \to 0
  \]
  of \(M/xM\) over \(R/xR\), of the same length \(e\). So
  \(\text{pd}_{R/xR}(M/xM) = \text{pd}_R(M)\). By Stacks
  \texttt{lemma-depth-drops-by-one}, both \(\text{depth}(R/xR) = \text{depth}(R) - 1\) and
  \(\text{depth}(M/xM) = \text{depth}(M) - 1\), and one verifies
  \(\text{depth}_R(M/xM) = \text{depth}_{R/xR}(M/xM)\) (regular sequences in \(\mathfrak{m}\)
  correspond to regular sequences in \(\mathfrak{m}/(x)\)). Apply the inductive hypothesis to
  \(M/xM\) over \(R/xR\):
  \[
    \text{pd}_{R/xR}(M/xM) + \text{depth}(M/xM) = \text{depth}(R/xR).
  \]
  Substituting yields
  \(\text{pd}_R(M) + (\text{depth}(M) - 1) = \text{depth}(R) - 1\), i.e.\
  \(\text{pd}_R(M) + \text{depth}(M) = \text{depth}(R)\), as required.
\end{proof}

\begin{lemma}[Auslander--Buchsbaum: inductive step \(\text{pd} = k+1\)]\label{mk-lem:auslander_buchsbaum_formula_succ_pd}
  \lean{RingTheory.auslander_buchsbaum_formula_succ_pd}
  \leanok
  \dcref{090V,custom}

  \uses{mk-lem:depth_drops_by_one, mk-lem:enat_ab_inductive_combine,
    mk-lem:projectiveDimension_ker_eq_of_surjection}
  \textit{Source: \cite{stacks-project}, tag 090V (proposition-Auslander-Buchsbaum),
  induction-on-\(\text{pd}\) re-organisation; cf.\ Matsumura, Theorem 19.1.}

  Let \((R,\mathfrak m)\) be a Noetherian local ring, let \(M\) be a
  nonzero finite \(R\)-module of projective dimension \(k+1\), and assume
  the Auslander--Buchsbaum formula for modules of projective dimension at
  most \(k\). Then
  \(\text{pd}_R(M) + \text{depth}(M) = \text{depth}(R)\).
\end{lemma}
\begin{proof}
  \leanok
  \uses{mk-lem:depth_drops_by_one, mk-lem:depth_ses_ineqs_of_surjection_finite_localRing, mk-lem:exists_ne_zero_ext_of_depth_eq, mk-lem:ext_comp_mk_zero_ofHom_eq_zero_of_entries_mem_annihilator, mk-lem:enat_ab_inductive_combine, mk-lem:projectiveDimension_ker_eq_of_surjection}
  Choose a minimal surjection \(f:R^n\twoheadrightarrow M\), and write
  \(K=\ker f\). The syzygy sequence
  \[
    0\longrightarrow K\longrightarrow R^n\xrightarrow{f}M\longrightarrow 0
  \]
  gives \(\operatorname{pd}_R(K)=k\) by
  \ref{mk-lem:projectiveDimension_ker_eq_of_surjection}.

  Suppose first that \(k>0\). By the inductive hypothesis,
  \(k+\operatorname{depth}(K)=\operatorname{depth}(R)\), so
  \(\operatorname{depth}(K)<\operatorname{depth}(R)\). The depth
  inequalities for the syzygy sequence then force
  \(\operatorname{depth}(K)=\operatorname{depth}(M)+1\). Substituting
  this equality into the inductive hypothesis gives the desired formula.

  It remains to treat \(k=0\). Then \(K\) is a nonzero finite free
  module, hence \(\operatorname{depth}(K)=\operatorname{depth}(R)=d\).
  Minimality of \(f\) means that every matrix entry of
  \(K\to R^n\) belongs to \(\mathfrak m=\operatorname{Ann}_R(R/\mathfrak m)\).
  Therefore the induced map
  \(\operatorname{Ext}^d_R(R/\mathfrak m,K)\to
  \operatorname{Ext}^d_R(R/\mathfrak m,R^n)\) is zero. The source is
  nonzero by the characterization of depth, so exactness of the long
  Ext sequence supplies a nonzero class in
  \(\operatorname{Ext}^{d-1}_R(R/\mathfrak m,M)\). Thus
  \(\operatorname{depth}(M)\leq d-1\), while the depth inequality for
  the same short exact sequence gives the reverse inequality. Hence
  \(\operatorname{depth}(M)=d-1\), and
  \(1+\operatorname{depth}(M)=\operatorname{depth}(R)\).
\end{proof}

\begin{lemma}

  [Depth drops by one under a regular element]
  \label{mk-lem:depth_drops_by_one}
  \dcref{00LL,custom}
  \uses{mk-def:depth, mk-lem:depth_via_ext}
  \textit{Source: \cite{stacks-project}, tag 00LL (\texttt{lemma-depth-drops-by-one}).}

  For a Noetherian local ring \((R, \mathfrak m)\), a nonzero finite
  \(R\)-module \(M\), and an \(M\)-regular element \(x \in \mathfrak m\),
  \[
    \text{depth}_{\mathfrak m}(M/xM) + 1 \;=\; \text{depth}_{\mathfrak m}(M).
  \]
\end{lemma}
\begin{proof}

  Reduces to the Ext characterisation
  (\ref{mk-lem:depth_via_ext}) on both sides; the short exact sequence
  \(0 \to M \xrightarrow{\,\cdot x\,} M \to M/xM \to 0\) gives a
  long exact sequence in \(\text{Ext}^*_R(\kappa, -)\) on which
  multiplication by \(x\) is zero (since \(x \in \text{Ann}_R \kappa\)),
  breaking the LES into short exact pieces
  \(0 \to \text{Ext}^i(\kappa, M) \to \text{Ext}^i(\kappa, M/xM) \to
  \text{Ext}^{i+1}(\kappa, M) \to 0\). The translation to the
  \(\mathbb N_\infty\) depth equality goes through
  \texttt{ENat.forall\_natCast\_le\_iff\_le}.
\end{proof}

\subsection{The syzygy step in the induction}
\label{mk-subsec:succ_pd_gap_sequence}

Let \(M\) have projective dimension \(k+1\), and choose a minimal surjection
\(f : R^n \twoheadrightarrow M\) with kernel \(K\). The short exact sequence
\[
  0 \longrightarrow K \longrightarrow R^n \xrightarrow{f} M \longrightarrow 0
\]
reduces the projective dimension by one: \(\operatorname{pd}_R(K)=k\).
The depth inequalities for this sequence compare
\(\operatorname{depth}(K)\), \(\operatorname{depth}(M)\), and
\(\operatorname{depth}(R^n)=\operatorname{depth}(R)\). For \(k>0\), the
inductive hypothesis and these inequalities give
\(\operatorname{depth}(K)=\operatorname{depth}(M)+1\). When \(k=0\), the
minimality of \(f\) makes the matrix of \(K\to R^n\) vanish after reduction
modulo the maximal ideal; the induced map on
\(\operatorname{Ext}^*_R(R/\mathfrak m,-)\) is therefore zero, and the same
depth equality follows from the long exact sequence. Substitution into the
inductive hypothesis yields
\[
  (k+1)+\operatorname{depth}(M)=\operatorname{depth}(R).
\]

\subsection{Minimal finite-free surjections}
\label{mk-subsec:ab_gap1_first_step}
\begin{lemma}[Per-syzygy minimal surjection from \(R^n\) for finite modules over a local ring]\label{mk-lem:exists_minimalSurjection_finite_localRing}
  \lean{RingTheory.Module.exists_minimalSurjection_finite_localRing}
  \leanok
  \dcref{00LK,custom}

  \uses{mk-def:depth}
  \textit{Source: \cite{stacks-project}, tag 00LK
    (\texttt{lemma-add-trivial-complex}).}
  Let \(R\) be a local ring (Noetherian-ness NOT required at this
  step), \(\mathfrak m \subseteq R\) the maximal ideal,
  \(\kappa = R/\mathfrak m\) the residue field, and \(M\) a finite
  \(R\)-module. Then there exist \(n \in \mathbb N\) and an
  \(R\)-linear surjection \(f : R^n \twoheadrightarrow M\) such that:
  \begin{enumerate}
    \item \(n = \dim_\kappa(\kappa \otimes_R M)\) (the minimal
      generator count of \(M\), via Nakayama);
    \item \(\ker f \subseteq \mathfrak m \cdot R^n\) (the
      ``minimality'' condition: every entry of every relation lives
      in the maximal ideal, so the chosen lift has no trivial
      summands).
  \end{enumerate}
\end{lemma}
\begin{proof}
  \leanok

Pick a \(\kappa\)-basis
  \(b : \mathrm{Fin}\,n \to \kappa \otimes_R M\) via
  \texttt{Module.finBasis} with
  \(n = \dim_\kappa(\kappa \otimes_R M)\). Surjectivity of the
  canonical map \(M \twoheadrightarrow \kappa \otimes_R M\)
  (composition of \texttt{TensorProduct.mk\_surjective} and
  \texttt{Ideal.Quotient.mk\_surjective}) gives lifts
  \(m_i \in M\) with \(1 \otimes m_i = b_i\). Define
  \(f : R^n \to M\) by \(f x := \sum_i x_i \cdot m_i\) via
  \texttt{(Pi.basisFun R (Fin n)).constr R m}.

  Surjectivity: \(\mathrm{range}\, f = \mathrm{span}_R(m_i)\) by
  \texttt{Module.Basis.constr\_range}, and
  \(\mathrm{span}_R(m_i) = M\) by Nakayama via
  \texttt{IsLocalRing.span\_eq\_top\_of\_tmul\_eq\_basis}.

  Kernel containment: for \(x \in \ker f\) the relation
  \(\sum_i x_i m_i = 0\) tensors to
  \(\sum_i (1 \otimes x_i \cdot m_i) = 0\); using
  \texttt{TensorProduct.tmul\_smul} and the linear independence of
  \(b\) (\texttt{Module.Basis.linearIndependent}) we obtain
  \(\mathrm{residue}_R(x_i) = 0\) for every \(i\), i.e.\
  \(x_i \in \mathfrak m\). Decomposing
  \(x = \sum_i x_i \cdot \mathrm{Pi.single}\,i\,1\) then yields
  \(x \in \mathfrak m \cdot R^n\) by
  \texttt{Submodule.smul\_mem\_smul} and
  \texttt{Submodule.sum\_mem}.
\end{proof}

The resulting short exact sequence is the syzygy sequence used in the
induction on projective dimension.

\subsection{Projective-dimension descent along a syzygy}
\label{mk-subsec:ab_gap1_haspdlt_pivot}

The following lemmas express the two directions of projective-dimension
control in the syzygy sequence and the corresponding lower bound on depth.
\begin{lemma}[Bridge: \(\texttt{projectiveDimension}\,M = n \Rightarrow \texttt{HasProjectiveDimensionLT}\,M\,(n+1)\)]\label{mk-lem:hasProjectiveDimensionLT_succ_of_projectiveDimension_eq}
  \lean{RingTheory.Module.hasProjectiveDimensionLT_succ_of_projectiveDimension_eq}
  \leanok
  \dcref{custom}

  \uses{mk-def:projective_dimension}
  \textit{Source: Mathlib's \(\texttt{CategoryTheory.projectiveDimension\_lt\_iff}\).}
  For a Noetherian local ring \(R\) and a finite \(R\)-module \(M\)
  with \(\texttt{Module.projectiveDimension}\,R\,M =
  ((n : \mathbb N) : \texttt{WithBot}\,\mathbb N_\infty)\), one has
  \(\texttt{HasProjectiveDimensionLT}\,(\texttt{ModuleCat.of}\,R\,M)\,
  (n+1)\). This is the strict-inequality formulation of the assumed
  projective-dimension equality.
\end{lemma}
\begin{lemma}[Syzygy descent of \(\texttt{HasProjectiveDimensionLT}\)]\label{mk-lem:hasProjectiveDimensionLT_ker_of_surjection}
  \lean{RingTheory.Module.hasProjectiveDimensionLT_ker_of_surjection}
  \leanok
  \dcref{custom}

  \uses{mk-lem:hasProjectiveDimensionLT_succ_of_projectiveDimension_eq}
  \textit{Source: Mathlib's
  \(\texttt{ShortComplex.ShortExact.hasProjectiveDimensionLT\_X\_1}\) +
  \(\texttt{ModuleCat.projective\_of\_free}\).}
  Let \(f : R^n \twoheadrightarrow M\) be a surjection of \(R\)-modules
  with \(\texttt{HasProjectiveDimensionLT}\,(\texttt{ModuleCat.of}\,R\,M)\,
  (k+2)\). Then
  \(\texttt{HasProjectiveDimensionLT}\,
   (\texttt{ModuleCat.of}\,R\,(\ker f))\,(k+1)\).
  Proof: build the SES \(0 \to \ker f \to R^n \to M \to 0\) via
  \(\texttt{LinearMap.shortExact\_shortComplexKer}\); discharge
  projectivity of \(R^n\) via \(\texttt{ModuleCat.projective\_of\_free}
   (\texttt{Pi.basisFun}\,R\,(\texttt{Fin}\,n))\); lift
  \(\texttt{HasProjectiveDimensionLT}\,1\) on \(R^n\) up to
  \(k+1\) via \(\texttt{hasProjectiveDimensionLT\_of\_ge}\); apply
  \(\texttt{ShortExact.hasProjectiveDimensionLT\_X\_1}\) to close.
\end{lemma}
\begin{lemma}[Ascent companion: \(\ker f\)-pd-control \(\Rightarrow M\)-pd-control]\label{mk-lem:hasProjectiveDimensionLT_succ_of_hasProjectiveDimensionLT_ker}
  \lean{RingTheory.Module.hasProjectiveDimensionLT_succ_of_hasProjectiveDimensionLT_ker}
  \leanok
  \dcref{custom}

  \uses{mk-lem:hasProjectiveDimensionLT_ker_of_surjection}
  Companion ascent direction: with the same SES,
  \(\texttt{HasProjectiveDimensionLT}\,(\ker f)\,(k+1)\) gives
  \(\texttt{HasProjectiveDimensionLT}\,M\,(k+2)\) via
  \(\texttt{ShortExact.hasProjectiveDimensionLT\_X\_3}\). Together with
  \ref{mk-lem:hasProjectiveDimensionLT_ker_of_surjection} this enables the
  contradiction argument that pins \(\texttt{pd}\,\ker f = k\)
  \emph{exactly} (rather than merely \(\leq k\)) under
  \(\texttt{pd}\,M = k+1\) exactly.
\end{lemma}
\begin{lemma}
[Depth lower bound on the kernel of a finite-free surjection]
  \label{mk-lem:depth_ker_ge_min_of_surjection_finite_localRing}
  \dcref{custom}
  \uses{mk-lem:depth_short_exact_sequence, mk-lem:depth_pi_const_eq_depth_of_nonempty}
  \textit{Source: \(\texttt{depth\_of\_short\_exact}\) part (3) +
  \(\texttt{depth\_pi\_const\_eq\_depth\_of\_nonempty}\).}
  Let \(R\) be a Noetherian local ring, \(M\) a nonzero finite
  \(R\)-module, \(n \geq 1\), and \(f : R^n \twoheadrightarrow M\) a
  surjection with \(\ker f\) nontrivial. Then
  \(\min(\texttt{depth}\,R,\, \texttt{depth}\,M + 1) \leq
   \texttt{depth}\,(\ker f)\).
  Proof: assemble the SES \(0 \to \ker f \to R^n \to M \to 0\) via
  \((\ker f).\texttt{subtype}\); rewrite
  \(\texttt{depth}\,(R^n) = \texttt{depth}\,R\) via
  \(\texttt{depth\_pi\_const\_eq\_depth\_of\_nonempty}\); apply
  \(\texttt{depth\_of\_short\_exact}\) part (3) for the lower bound.
\end{lemma}

Together, the preceding lemmas identify the projective dimension of the
syzygy and supply the depth comparison used in
\ref{mk-lem:auslander_buchsbaum_formula_succ_pd}. The base case is completed
by the Ext-annihilation lemmas recorded below.
\section{Auxiliary homological lemmas}
\label{mk-sec:ab_substrate_helpers}

This section records the elementary Ext, depth, and regular-sequence lemmas
used in the proofs above.

\subsection{Ext-annihilation and depth-bound substrate}
\label{mk-subsec:ab_ext_depth_substrate}

\begin{lemma}[Annihilator kills \(\text{Ext}\)]\label{mk-lem:ext_smul_eq_zero_of_mem_annihilator}
  \lean{RingTheory.Module.ext_smul_eq_zero_of_mem_annihilator}
  \leanok
  \dcref{custom}

  \uses{mk-def:depth}
  Let \(R\) be a commutative ring, let \(N, M\) be \(R\)-modules, fix \(i \in \mathbb N\),
  and let \(e \in \text{Ext}^i_R(N, M)\). If \(x \in R\) annihilates \(N\) (i.e.\
  \(x \in \text{Ann}_R(N)\)), then the \(R\)-action \(x \cdot e = 0\). Concretely,
  \(x \cdot e = (\text{mk}_0(x \cdot \mathbf 1_N)) \circ e\), and \(x \cdot \mathbf 1_N = 0\)
  since \(x\) annihilates \(N\).
\end{lemma}
\begin{proof}
  \leanok

The scalar action of \(x\) on \(\operatorname{Ext}^i_R(N,M)\) is induced by
  precomposition with multiplication by \(x\) on \(N\). Since \(xN=0\),
  that endomorphism is zero, and functoriality sends it to the zero map.
\end{proof}

\begin{lemma}[\(\text{Ext}\)-vanishing below the depth]\label{mk-lem:ext_vanish_of_natCast_lt_depth}
  \lean{RingTheory.Module.ext_vanish_of_natCast_lt_depth}
  \leanok
  \dcref{custom}

  \uses{mk-def:depth, mk-lem:depth_via_ext}
  Let \((R, \mathfrak m)\) be a Noetherian local ring with residue field \(\kappa\),
  and let \(M\) be a nonzero finitely generated \(R\)-module. If
  \((i : \mathbb N_\infty) < \text{depth}(M)\), then every element of
  \(\text{Ext}^i_R(\kappa, M)\) is zero. This is the index-wise repackaging of
  \ref{mk-lem:depth_via_ext} convenient for the long-exact-sequence chases.
\end{lemma}
\begin{proof}
  \leanok

If \(\operatorname{Ext}^i_R(\kappa,M)\) were nonzero, the
  characterization in \ref{mk-lem:depth_via_ext} would give
  \(\operatorname{depth}(M)\leq i\), contrary to the hypothesis.
\end{proof}

\begin{lemma}[\(\mathbb N_\infty\) successor bridge]\label{mk-lem:natCast_add_one_le_of_le_sub_one}
  \lean{RingTheory.Module.natCast_add_one_le_of_le_sub_one}
  \leanok
  \dcref{custom}

  For \(d \in \mathbb N_\infty\) and \(a \in \mathbb N\) with \(1 \leq a\): if
  \((a : \mathbb N_\infty) \leq d - 1\) then \((a + 1 : \mathbb N_\infty) \leq d\). A
  truncated-subtraction arithmetic lemma used for the \(\text{depth}(N') - 1\)
  shift in the second inequality of \ref{mk-lem:depth_short_exact_sequence}.
\end{lemma}
\begin{proof}
  \leanok

If \(d=\infty\), the conclusion is immediate. Otherwise write
  \(d=m\). Truncated subtraction agrees with ordinary subtraction under
  \(1\leq a\leq m-1\), so \(a+1\leq m=d\).
\end{proof}

\begin{lemma}[Depth is a linear-equivalence invariant]\label{mk-lem:depth_eq_of_linearEquiv}
  \lean{RingTheory.Module.depth_eq_of_linearEquiv}
  \leanok
  \dcref{custom}

  \uses{mk-def:depth}
  Let \(R\) be a commutative ring, \(I \subseteq R\) an ideal, and \(M, M'\) two
  \(R\)-modules related by an \(R\)-linear equivalence \(e : M \xrightarrow{\sim} M'\).
  Then \(\text{depth}_I(M) = \text{depth}_I(M')\): both the side condition
  \(I \cdot M = M\) and the regular-sequence supremum are preserved under \(e\).
\end{lemma}
\begin{proof}
  \leanok

The equivalence carries \(I M\) onto \(I M'\), so the exceptional
  case \(IM=M\) is preserved. It also carries each regular sequence on
  \(M\) to a regular sequence of the same length on \(M'\), and conversely
  via the inverse equivalence. The two suprema are therefore equal.
\end{proof}

\subsection{Depth of a constant product module}
\label{mk-subsec:ab_depth_pi_substrate}

\begin{lemma}[Ideal action on a constant product]\label{mk-lem:ideal_smul_top_pi_const}
  \lean{RingTheory.Module.ideal_smul_top_pi_const}
  \leanok
  \dcref{custom}

  For a commutative ring \(R\), an ideal \(I\), a finite index type \(\iota\), and an
  \(R\)-module \(M\), the submodule \(I \cdot \top_{\iota \to M}\) equals the product
  submodule \(\prod_{\iota} (I \cdot \top_M)\) of fibrewise ideal actions.
\end{lemma}
\begin{proof}
  \leanok

For one inclusion, take coordinates in a finite sum of elements
  \(a m\) with \(a\in I\); every coordinate lies in \(IM\). Conversely,
  decompose a tuple into its finitely many coordinate vectors. Each vector
  lies in the ideal action because its unique nonzero coordinate does.
\end{proof}

\begin{lemma}[Triviality of the ideal action on a constant product]\label{mk-lem:ideal_smul_top_pi_const_eq_top_iff}
  \lean{RingTheory.Module.ideal_smul_top_pi_const_eq_top_iff}
  \leanok
  \dcref{custom}

  \uses{mk-lem:ideal_smul_top_pi_const}
  For a nonempty finite index type \(\iota\): \(I \cdot \top_{\iota \to M} = \top\) if
  and only if \(I \cdot \top_M = \top\). The forward direction reads off the fibre
  witness from a \(\text{Pi.single}\) projection; the backward direction lifts a
  fibre witness through \(\text{Pi.single}\).
\end{lemma}
\begin{proof}
  \leanok

Apply \ref{mk-lem:ideal_smul_top_pi_const}. If \(IM=M\), the product
  equality follows coordinatewise. Conversely, evaluate the product
  equality at any index, which exists by nonemptiness, to obtain \(IM=M\).
\end{proof}

\begin{definition}[Quotient of a constant product by a scalar]\label{mk-def:quotSMulTopPiConstLinearEquiv}
  \lean{RingTheory.Module.quotSMulTopPiConstLinearEquiv}
  \leanok
  \dcref{custom}

  \uses{mk-lem:ideal_smul_top_pi_const}
  For a commutative ring \(R\), a finite index type \(\iota\), a scalar \(r \in R\), and
  an \(R\)-module \(M\), the canonical \(R\)-linear equivalence
  \[
    (\iota \to M)\big/ r \cdot \top \;\xrightarrow{\ \sim\ }\; \iota \to (M / r \cdot \top),
  \]
  obtained by rewriting \(r \cdot \top = \text{span}\{r\} \cdot \top\) and applying the
  quotient-of-a-product identification.
\end{definition}

\begin{lemma}[Regularity on a constant product]\label{mk-lem:isRegular_pi_const_iff_of_nonempty}
  \lean{RingTheory.Module.isRegular_pi_const_iff_of_nonempty}
  \leanok
  \dcref{custom}

  \uses{mk-def:quotSMulTopPiConstLinearEquiv}
  For a nonempty finite index type \(\iota\) and a list \(rs\) of elements of \(R\), the
  sequence \(rs\) is \((\iota \to M)\)-regular if and only if it is \(M\)-regular. The
  proof inducts on \(rs\): the empty case reduces to
  \(\text{Nontrivial}(\iota \to M) \leftrightarrow \text{Nontrivial}(M)\), and the cons
  case combines pointwise \(\text{SMul}\)-regularity with
  \ref{mk-def:quotSMulTopPiConstLinearEquiv} to feed the inductive hypothesis.
\end{lemma}
\begin{proof}
  \leanok

Induct on the sequence. Nontriviality of a nonempty finite product
  is equivalent to nontriviality of one factor. For the inductive step,
  scalar regularity is coordinatewise, and
  \ref{mk-def:quotSMulTopPiConstLinearEquiv} identifies the quotient product
  with the product of the quotients; apply the inductive hypothesis.
\end{proof}

\begin{lemma}[Depth of a constant product]\label{mk-lem:depth_pi_const_eq_depth_of_nonempty}
  \lean{RingTheory.Module.depth_pi_const_eq_depth_of_nonempty}
  \leanok
  \dcref{custom}

  \uses{mk-def:depth, mk-lem:ideal_smul_top_pi_const_eq_top_iff, mk-lem:isRegular_pi_const_iff_of_nonempty}
  For a commutative ring \(R\), an ideal \(I\), an \(R\)-module \(M\), and a nonempty
  finite index type \(\iota\), one has \(\text{depth}_I(\iota \to M) = \text{depth}_I(M)\).
  In particular a finite free module \(M \cong R^k\) has \(\text{depth}(M) = \text{depth}(R)\),
  which closes the \(\text{pd}_R(M) = 0\) base case of \ref{mk-thm:auslander_buchsbaum}.
\end{lemma}
\begin{proof}
  \leanok

The exceptional condition in the definition of depth is equivalent
  by \ref{mk-lem:ideal_smul_top_pi_const_eq_top_iff}. The regular sequences,
  with their lengths, agree by
  \ref{mk-lem:isRegular_pi_const_iff_of_nonempty}. Hence their suprema agree.
\end{proof}

\subsection{Surjection and projective-dimension substrate for the inductive step}
\label{mk-subsec:ab_surjection_substrate}

\begin{lemma}
[A regular element exists when depth is positive]
  \label{mk-lem:exists_isSMulRegular_of_one_le_depth}
  \dcref{custom}
  \uses{mk-def:depth}
  For a Noetherian local ring \((R, \mathfrak m)\) and a nonzero finitely generated
  \(R\)-module \(M\), if \(1 \leq \text{depth}(M)\) then there exists \(x \in \mathfrak m\)
  that is \(M\)-regular. (Nakayama rules out the trivialisation
  \(\mathfrak m \cdot M = M\), so the depth supremum exhibits a length-\(1\) witness.)
\end{lemma}
\begin{proof}

  Nakayama gives \(\mathfrak mM\neq M\). Since depth is at least one,
  its defining supremum contains a regular sequence of positive length.
  The first element of that sequence lies in \(\mathfrak m\) and is
  \(M\)-regular.
\end{proof}

\begin{lemma}[Kernel-SES depth inequalities for a finite free surjection]\label{mk-lem:depth_ses_ineqs_of_surjection_finite_localRing}
  \lean{RingTheory.Module.depth_ses_ineqs_of_surjection_finite_localRing}
  \leanok
  \dcref{custom}

  \uses{mk-lem:depth_short_exact_sequence, mk-lem:depth_pi_const_eq_depth_of_nonempty}
  Let \(R\) be a Noetherian local ring, \(M\) a nonzero finite \(R\)-module, \(n \geq 1\),
  and \(f : R^n \twoheadrightarrow M\) a surjection with \(\ker f\) nontrivial. Then
  \[
    \min\bigl(\text{depth}(R),\, \text{depth}(\ker f) - 1\bigr) \leq \text{depth}(M)
    \quad\text{and}\quad
    \min\bigl(\text{depth}(R),\, \text{depth}(M) + 1\bigr) \leq \text{depth}(\ker f).
  \]
  These are parts (2) and (3) of \ref{mk-lem:depth_short_exact_sequence} for the SES
  \(0 \to \ker f \to R^n \to M \to 0\), after identifying \(\text{depth}(R^n) = \text{depth}(R)\).
\end{lemma}
\begin{proof}
  \leanok

Apply parts (2) and (3) of
  \ref{mk-lem:depth_short_exact_sequence} to
  \(0\to\ker f\to R^n\to M\to0\), and replace
  \(\operatorname{depth}(R^n)\) by \(\operatorname{depth}(R)\) using
  \ref{mk-lem:depth_pi_const_eq_depth_of_nonempty}.
\end{proof}

\begin{lemma}[A nonzero \(\text{Ext}\) class at the depth index]\label{mk-lem:exists_ne_zero_ext_of_depth_eq}
  \lean{RingTheory.Module.exists_ne_zero_ext_of_depth_eq}
  \leanok
  \dcref{custom}

  \uses{mk-lem:depth_via_ext}
  For a nonzero finite \(R\)-module \(M\) over a Noetherian local ring \((R, \mathfrak m)\)
  with \(\text{depth}(M) = D\) finite, there is a nonzero element of
  \(\text{Ext}^D_R(\kappa, M)\). (All \(\text{Ext}^i\) vanish for \(i < D\), and
  vanishing fails at index \(D\).) This exhibits the nonzero class
  \(\text{Ext}^{\text{depth}(R)}(\kappa, R^k) \neq 0\) in the base case of the formula.
\end{lemma}
\begin{proof}
  \leanok

By \ref{mk-lem:depth_via_ext}, \(D\) is the least index at which
  \(\operatorname{Ext}^i_R(\kappa,M)\) does not vanish. Thus the group at
  index \(D\) is nonzero, and any nonzero group contains the required class.
\end{proof}

\begin{lemma}[Exact projective dimension of a syzygy]\label{mk-lem:projectiveDimension_ker_eq_of_surjection}
  \lean{RingTheory.projectiveDimension_ker_eq_of_surjection}
  \leanok
  \dcref{custom}

  \uses{mk-def:projective_dimension, mk-lem:hasProjectiveDimensionLT_succ_of_projectiveDimension_eq, mk-lem:hasProjectiveDimensionLT_ker_of_surjection, mk-lem:hasProjectiveDimensionLT_succ_of_hasProjectiveDimensionLT_ker}
  For a surjection \(f : R^n \twoheadrightarrow M\) with \(\text{pd}_R(M) = k + 2\), the
  kernel satisfies \(\text{pd}_R(\ker f) = k + 1\) exactly. The upper bound is the
  syzygy descent of \ref{mk-lem:hasProjectiveDimensionLT_ker_of_surjection}; the lower
  bound is the contrapositive of the ascent
  \ref{mk-lem:hasProjectiveDimensionLT_succ_of_hasProjectiveDimensionLT_ker}. This is
  the exact-\(\text{pd}\) input that licenses the inductive hypothesis on \(\ker f\).
\end{lemma}
\begin{proof}
  \leanok

Syzygy descent gives \(\operatorname{pd}_R(\ker f)\leq k+1\).
  If it were at most \(k\), the ascent statement applied to the same short
  exact sequence would give \(\operatorname{pd}_R(M)\leq k+1\), contradicting
  the assumed value \(k+2\).
\end{proof}

\begin{lemma}[\(\mathbb N_\infty\) combine for the inductive step]\label{mk-lem:enat_ab_inductive_combine}
  \lean{RingTheory.enat_ab_inductive_combine}
  \leanok
  \dcref{custom}

  A purely arithmetic statement in \(\mathbb N_\infty\): given \(j + d_K = d\),
  \(\min(d, d_K - 1) \leq d_M\), \(\min(d, d_M + 1) \leq d_K\), and \(1 \leq j\), one
  concludes \((j + 1) + d_M = d\). This packages the inductive hypothesis
  \(j + \text{depth}(\ker f) = \text{depth}(R)\) with the two kernel-SES depth
  inequalities into the next instance of the Auslander--Buchsbaum equation.
\end{lemma}
\begin{proof}
  \leanok

The equality \(j+d_K=d\) and \(1\leq j\) exclude \(d_K=\infty\)
  unless \(d=\infty\), where the conclusion is immediate. In the finite
  case the two minimum inequalities give
  \(d_M+1\leq d_K\leq d_M+1\); hence \(d_K=d_M+1\), and substitution yields
  \((j+1)+d_M=d\).
\end{proof}

\subsection{Matrix-collapse on \texorpdfstring{\(\text{Ext}\)}{Ext}}
\label{mk-subsec:ab_matrix_collapse_substrate}

\begin{definition}[Elementary linear map]\label{mk-def:elemMap}
  \lean{RingTheory.Module.elemMap}
  \leanok
  \dcref{custom}

  For a commutative ring \(R\) and indices \(i \in \{1, \dots, n\}\),
  \(j \in \{1, \dots, m\}\), the \emph{elementary map} \(E_{i,j} : R^m \to R^n\) is the
  \(R\)-linear map sending the standard basis vector \(\text{Pi.single}\,j\,1\) to
  \(\text{Pi.single}\,i\,1\) and every other standard basis vector to \(0\).
\end{definition}

\begin{lemma}[Evaluation of the elementary map]\label{mk-lem:elemMap_apply}
  \lean{RingTheory.Module.elemMap_apply}
  \leanok
  \dcref{custom}

  \uses{mk-def:elemMap}
  For \(x \in R^m\) one has \(E_{i,j}(x) = \text{Pi.single}\,i\,(x_j)\): the elementary
  map reads off the \(j\)-th coordinate and places it in the \(i\)-th slot.
\end{lemma}
\begin{proof}
  \leanok

Expand \(x\) in the standard basis. By linearity, \(E_{i,j}\)
  kills every basis term except the \(j\)-th and sends that coefficient to
  the \(i\)-th coordinate, giving the displayed formula.
\end{proof}

\begin{lemma}[Matrix decomposition into elementary maps]\label{mk-lem:linearMap_finFunR_matrix_decomp}
  \lean{RingTheory.Module.linearMap_finFunR_matrix_decomp}
  \leanok
  \dcref{custom}

  \uses{mk-def:elemMap, mk-lem:elemMap_apply}
  Every \(R\)-linear map \(A : R^m \to R^n\) decomposes as
  \(A = \sum_{i,j} A_{i,j} \cdot E_{i,j}\), where \(A_{i,j} = (A\,(\text{Pi.single}\,j\,1))_i\)
  is the corresponding matrix entry.
\end{lemma}
\begin{proof}
  \leanok

Both sides are linear maps, so it suffices to evaluate them on
  each standard basis vector. The \(j\)-th column of the sum is
  \(\sum_i A_{i,j}e_i=A(e_j)\), proving equality.
\end{proof}

\begin{lemma}[Matrix-collapse on \(\text{Ext}\)]\label{mk-lem:ext_comp_mk_zero_ofHom_eq_zero_of_entries_mem_annihilator}
  \lean{RingTheory.Module.ext_comp_mk₀_ofHom_eq_zero_of_entries_mem_annihilator}
  \leanok
  \dcref{custom}

  \uses{mk-lem:linearMap_finFunR_matrix_decomp, mk-lem:ext_smul_eq_zero_of_mem_annihilator}
  Let \(A : R^m \to R^n\) be an \(R\)-linear map every matrix entry of which lies in
  \(\text{Ann}_R(N)\). Then the postcomposition map
  \(\text{Ext}^p_R(N, R^m) \to \text{Ext}^p_R(N, R^n)\) induced by \(A\) is identically
  zero. The proof decomposes \(A = \sum_{i,j} A_{i,j} \cdot E_{i,j}\)
  (\ref{mk-lem:linearMap_finFunR_matrix_decomp}) and kills each scalar summand via
  \ref{mk-lem:ext_smul_eq_zero_of_mem_annihilator}.
\end{lemma}
\begin{proof}
  \leanok

Use the elementary-map decomposition of \(A\). The induced Ext map
  is additive and linear in each matrix coefficient. Every coefficient
  annihilates \(N\), so \ref{mk-lem:ext_smul_eq_zero_of_mem_annihilator}
  kills every summand; their sum is zero.
\end{proof}

\subsection{Regular-sequence and domain substrate for regular local rings}
\label{mk-subsec:ab_regular_domain_substrate}

\begin{lemma}[Regular sequences are bounded by the Krull dimension]\label{mk-lem:length_le_ringKrullDim_of_isRegular}
  \lean{RingTheory.CohenMacaulay.length_le_ringKrullDim_of_isRegular}
  \leanok
  \dcref{custom}

  For a Noetherian local ring \(R\), every \(R\)-regular sequence \(rs\) has length at
  most \(\dim(R)\). This is the upper-bound half of \ref{mk-cor:regular_cohen_macaulay},
  obtained from the dimension identity \(\dim(R / (rs)) + |rs| = \dim(R)\) together with
  \(\dim(R / (rs)) \geq 0\).
\end{lemma}
\begin{proof}
  \leanok

Quotient successively by the elements of the regular sequence.
  Each regular element lowers the Krull dimension by one, while the final
  quotient has nonnegative dimension. Thus
  \(\dim(R)\geq |rs|\).
\end{proof}

\begin{lemma}[Cotangent image of \(x \in \mathfrak m \setminus \mathfrak m^2\)]\label{mk-lem:toCotangent_ne_zero_of_not_mem_sq}
  \lean{RingTheory.CohenMacaulay.toCotangent_ne_zero_of_not_mem_sq}
  \leanok
  \dcref{custom}

  For a local ring \((R, \mathfrak m)\) and \(x \in \mathfrak m\) with \(x \notin \mathfrak m^2\),
  the image of \(x\) in the cotangent space \(\mathfrak m / \mathfrak m^2\) is nonzero. This is
  the positivity input for the cotangent dimension-drop.
\end{lemma}
\begin{proof}
  \leanok

The kernel of the quotient map
  \(\mathfrak m\to\mathfrak m/\mathfrak m^2\) is exactly
  \(\mathfrak m^2\). Hence the class of \(x\) vanishes precisely when
  \(x\in\mathfrak m^2\).
\end{proof}

\begin{lemma}[Cotangent dimension drops by one modulo \(x\)]\label{mk-lem:finrank_cotangentSpace_quot_span_singleton_succ}
  \lean{RingTheory.CohenMacaulay.finrank_cotangentSpace_quot_span_singleton_succ}
  \leanok
  \dcref{custom}

  \uses{mk-lem:toCotangent_ne_zero_of_not_mem_sq}
  For a Noetherian local ring \((R, \mathfrak m)\) and \(x \in \mathfrak m \setminus \mathfrak m^2\),
  the cotangent space of \(R / (x)\) has dimension one less than that of \(R\):
  \[
    \dim_{\kappa'}\!\bigl(\text{CotangentSpace}(R/(x))\bigr) + 1
      = \dim_{\kappa}\!\bigl(\text{CotangentSpace}(R)\bigr),
  \]
  where \(\kappa, \kappa'\) are the residue fields of \(R\) and \(R/(x)\). The kernel of
  the induced surjection on cotangent spaces is the line spanned by the image of \(x\),
  which is one-dimensional by \ref{mk-lem:toCotangent_ne_zero_of_not_mem_sq}.
\end{lemma}
\begin{proof}
  \leanok

The quotient map induces a surjection from the cotangent space of
  \(R\) onto that of \(R/(x)\). Its kernel is the line generated by the
  nonzero class of \(x\), so rank-nullity lowers the dimension by exactly one.
\end{proof}

\begin{lemma}[Nakayama witness for a positive cotangent dimension]\label{mk-lem:exists_notMemSq_of_spanFinrank_pos}
  \lean{RingTheory.CohenMacaulay.exists_notMemSq_of_spanFinrank_pos}
  \leanok
  \dcref{custom}

  For a Noetherian local ring \((R, \mathfrak m)\) with positive minimal-generator count
  of \(\mathfrak m\) (i.e.\ \(\text{spanFinrank}(\mathfrak m) \geq 1\)), there exists
  \(x \in \mathfrak m\) with \(x \notin \mathfrak m^2\). By Nakayama, \(\mathfrak m \subseteq \mathfrak m^2\)
  would force \(\mathfrak m = 0\), contradicting positivity.
\end{lemma}
\begin{proof}
  \leanok

If every element of \(\mathfrak m\) lay in \(\mathfrak m^2\), then
  \(\mathfrak m=\mathfrak m^2\). Nakayama's lemma would force
  \(\mathfrak m=0\), making the cotangent dimension zero, a contradiction.
\end{proof}

\begin{lemma}[Zero cotangent dimension forces a field]\label{mk-lem:isDomain_of_isLocalRing_of_spanFinrank_maximalIdeal_eq_zero}
  \lean{RingTheory.CohenMacaulay.isDomain_of_isLocalRing_of_spanFinrank_maximalIdeal_eq_zero}
  \leanok
  \dcref{custom}

  For a Noetherian local ring \(R\) with \(\text{spanFinrank}(\mathfrak m) = 0\), the maximal
  ideal collapses to \(0\), so \(R\) is a field and hence an integral domain. This is the
  base case of \ref{mk-lem:isDomain_of_regularLocal}.
\end{lemma}
\begin{proof}
  \leanok

Vanishing cotangent dimension gives
  \(\mathfrak m/\mathfrak m^2=0\), hence
  \(\mathfrak m=\mathfrak m^2\). Nakayama yields \(\mathfrak m=0\).
  A local ring whose maximal ideal is zero is a field, and therefore a domain.
\end{proof}

\begin{lemma}[Regularity descends to \(R/(x)\)]\label{mk-lem:regularLocal_quotient_isRegularLocal_of_notMemSq}
  \lean{RingTheory.CohenMacaulay.regularLocal_quotient_isRegularLocal_of_notMemSq}
  \leanok
  \dcref{custom}

  \uses{mk-lem:finrank_cotangentSpace_quot_span_singleton_succ}
  For a regular Noetherian local ring \(R\) with \(\text{spanFinrank}(\mathfrak m) = k + 1\)
  and \(x \in \mathfrak m \setminus \mathfrak m^2\), the quotient \(R / (x)\) is again a regular
  local ring, with \(\text{spanFinrank}(\mathfrak m') = k\). The cotangent dimension drop
  (\ref{mk-lem:finrank_cotangentSpace_quot_span_singleton_succ}) plus a Krull-height bound
  give the matching dimension equality, so \(R/(x)\) is regular. Notably this avoids
  assuming \(x\) is a non-zero-divisor, so it does not depend on
  \ref{mk-lem:isDomain_of_regularLocal}.
\end{lemma}
\begin{proof}
  \leanok

By \ref{mk-lem:finrank_cotangentSpace_quot_span_singleton_succ},
  the embedding dimension of \(R/(x)\) is \(k\). The principal ideal
  theorem gives \(\dim(R/(x))\geq\dim(R)-1=k\), while Krull dimension is
  bounded above by embedding dimension. Equality follows, which is the
  regularity criterion for the local quotient.
\end{proof}

\begin{lemma}[Members of a minimal prime are zero-divisors]\label{mk-lem:exists_ne_zero_mul_eq_zero_of_mem_minimalPrimes}
  \lean{RingTheory.CohenMacaulay.exists_ne_zero_mul_eq_zero_of_mem_minimalPrimes}
  \leanok
  \dcref{custom}

  For a commutative ring \(R\) and a minimal prime \(\mathfrak p\) of \(R\), every
  \(x \in \mathfrak p\) is a zero-divisor: there exists \(y \neq 0\) with \(x y = 0\). This
  follows from the disjointness of a minimal prime from the non-zero-divisors.
\end{lemma}
\begin{proof}
  \leanok

A minimal prime is disjoint from the set of non-zero-divisors. Hence
  \(x\) is not a non-zero-divisor, so multiplication by \(x\) is not
  injective. Thus there is \(y\neq0\) with \(xy=0\).
\end{proof}

\begin{lemma}[\((x)\) is not a minimal prime in the inductive step]\label{mk-lem:notMem_minimalPrimes_of_regularLocal_succ}
  \lean{RingTheory.CohenMacaulay.notMem_minimalPrimes_of_regularLocal_succ}
  \leanok
  \dcref{custom}

  \uses{mk-lem:exists_ne_zero_mul_eq_zero_of_mem_minimalPrimes, mk-lem:regularLocal_quotient_isRegularLocal_of_notMemSq}
  Let \(R\) be a regular Noetherian local ring with \(\text{spanFinrank}(\mathfrak m) = k + 1\)
  and \(x \in \mathfrak m \setminus \mathfrak m^2\), and suppose every \(k\)-dimensional regular
  local ring is a domain. Then \(\text{span}\{x\}\) is \emph{not} a minimal prime of \(R\).
  Combined with the regular descent
  (\ref{mk-lem:regularLocal_quotient_isRegularLocal_of_notMemSq}), this is the residual
  case of the regular-local domain argument.
\end{lemma}
\begin{proof}
  \leanok

Suppose that \((x)\) is minimal. Since a Noetherian ring has finitely many
  minimal primes, prime avoidance supplies
  \(x'\in\mathfrak m\) lying neither in \(\mathfrak m^2\) nor in any minimal
  prime: indeed, \(\mathfrak m\not\subseteq\mathfrak m^2\) by positive
  cotangent dimension, and \(\mathfrak m\) cannot lie in a minimal prime
  without forcing \(\dim(R)=0\).

  The quotient \(R/(x')\) is regular of dimension \(k\), hence a domain by
  the inductive hypothesis, so \((x')\) is prime. Choose a minimal prime
  \(\mathfrak p\subseteq(x')\). Since \(x'\notin\mathfrak p\), every
  \(y\in\mathfrak p\), written \(y=x'z\), has \(z\in\mathfrak p\) by
  primality. Thus \(\mathfrak p\subseteq\mathfrak m\mathfrak p\), and the
  reverse inclusion is automatic. The ideal \(\mathfrak p\) is finitely
  generated, so Nakayama gives \(\mathfrak p=0\). Consequently \(R\) is a
  domain. Its only minimal prime is then \(0\), so the assumed minimality of
  \((x)\) gives \(x=0\), contradicting \(x\notin\mathfrak m^2\).
\end{proof}

\begin{lemma}[Regular local rings are domains]\label{mk-lem:isDomain_of_regularLocal}
  \lean{RingTheory.CohenMacaulay.isDomain_of_regularLocal}
  \leanok
  \dcref{custom}

  \uses{mk-lem:isDomain_of_isLocalRing_of_spanFinrank_maximalIdeal_eq_zero, mk-lem:exists_notMemSq_of_spanFinrank_pos, mk-lem:regularLocal_quotient_isRegularLocal_of_notMemSq, mk-lem:notMem_minimalPrimes_of_regularLocal_succ}
  Every regular Noetherian local ring \(R\) is an integral domain. The proof is a strong
  induction on \(\text{spanFinrank}(\mathfrak m)\): the base case is the field case
  (\ref{mk-lem:isDomain_of_isLocalRing_of_spanFinrank_maximalIdeal_eq_zero}); the inductive
  step picks \(x \in \mathfrak m \setminus \mathfrak m^2\)
  (\ref{mk-lem:exists_notMemSq_of_spanFinrank_pos}), uses that \(R/(x)\) is regular local of
  smaller dimension (\ref{mk-lem:regularLocal_quotient_isRegularLocal_of_notMemSq}) and a
  domain by induction, and rules out the obstruction case via
  \ref{mk-lem:notMem_minimalPrimes_of_regularLocal_succ}.
\end{lemma}
\begin{proof}
  \leanok

Induct on the cotangent dimension. The zero-dimensional case is
  \ref{mk-lem:isDomain_of_isLocalRing_of_spanFinrank_maximalIdeal_eq_zero}.
  In positive dimension choose \(x\in\mathfrak m\setminus\mathfrak m^2\).
  The quotient is regular of smaller dimension and hence a domain by
  induction. The only possible obstruction to \(R\) being a domain would
  put \((x)\) among the minimal primes, contrary to
  \ref{mk-lem:notMem_minimalPrimes_of_regularLocal_succ}.
\end{proof}

\begin{lemma}[A regular element outside \(\mathfrak m^2\)]\label{mk-lem:exists_isSMulRegular_notMemSq_of_regularLocal_succ}
  \lean{RingTheory.CohenMacaulay.exists_isSMulRegular_notMemSq_of_regularLocal_succ}
  \leanok
  \dcref{custom}

  \uses{mk-lem:exists_notMemSq_of_spanFinrank_pos, mk-lem:isDomain_of_regularLocal}
  For a regular local ring \(R\) of dimension \(k + 1\), there exists
  \(x \in \mathfrak m \setminus \mathfrak m^2\) that is moreover an \(R\)-regular element. The
  Nakayama witness (\ref{mk-lem:exists_notMemSq_of_spanFinrank_pos}) is nonzero, and in a
  domain (\ref{mk-lem:isDomain_of_regularLocal}) every nonzero element is a
  non-zero-divisor.
\end{lemma}
\begin{proof}
  \leanok

Choose \(x\in\mathfrak m\setminus\mathfrak m^2\) by
  \ref{mk-lem:exists_notMemSq_of_spanFinrank_pos}. It is nonzero, and
  \ref{mk-lem:isDomain_of_regularLocal} makes multiplication by \(x\)
  injective. Thus \(x\) is regular.
\end{proof}

\begin{lemma}[Regular element with regular quotient]\label{mk-lem:exists_isSMulRegular_quotient_isRegularLocal_succ}
  \lean{RingTheory.CohenMacaulay.exists_isSMulRegular_quotient_isRegularLocal_succ}
  \leanok
  \dcref{custom}

  \uses{mk-lem:exists_isSMulRegular_notMemSq_of_regularLocal_succ, mk-lem:finrank_cotangentSpace_quot_span_singleton_succ}
  For a regular local ring \(R\) of dimension \(k + 1\), there exists \(x \in \mathfrak m\)
  that is \(R\)-regular and such that \(R / (x)\) is again a regular local ring of
  dimension \(k\) (its maximal ideal has \(\text{spanFinrank} = k\)). This consolidates the
  ``pick an \(R\)-regular system-of-parameters element and drop dimension'' step.
\end{lemma}
\begin{proof}
  \leanok

Take the element supplied by
  \ref{mk-lem:exists_isSMulRegular_notMemSq_of_regularLocal_succ}.
  It is regular, and its quotient is regular of dimension \(k\) by
  \ref{mk-lem:regularLocal_quotient_isRegularLocal_of_notMemSq}.
\end{proof}

\begin{lemma}[Inductive step for the regular sequence]\label{mk-lem:regularLocal_inductive_step}
  \lean{RingTheory.CohenMacaulay.regularLocal_inductive_step}
  \leanok
  \dcref{custom}

  \uses{mk-lem:exists_isSMulRegular_quotient_isRegularLocal_succ}
  Let \(R\) be a regular local ring of dimension \(k + 1\), and suppose that every regular
  local ring of dimension \(k\) admits a length-\(k\) regular sequence inside its maximal
  ideal. Then \(R\) admits a length-\((k+1)\) regular sequence inside \(\mathfrak m\): extract
  the regular element \(x\) with regular quotient
  (\ref{mk-lem:exists_isSMulRegular_quotient_isRegularLocal_succ}), lift a length-\(k\)
  regular sequence of \(R/(x)\) back to \(R\), and prepend \(x\).
\end{lemma}
\begin{proof}
  \leanok

Choose a regular \(x\) whose quotient is regular of dimension
  \(k\). By the inductive hypothesis, \(R/(x)\) has a regular sequence
  \(\bar x_1,\ldots,\bar x_k\). Lift these elements to \(\mathfrak m\);
  the definition of regularity in successive quotients shows that
  \(x,x_1,\ldots,x_k\) is regular on \(R\).
\end{proof}

\begin{lemma}[Regular local rings have a full regular sequence]\label{mk-lem:exists_isRegular_of_regularLocal}
  \lean{RingTheory.CohenMacaulay.exists_isRegular_of_regularLocal}
  \leanok
  \dcref{custom}

  \uses{mk-lem:regularLocal_inductive_step}
  For a regular Noetherian local ring \(R\), there is an \(R\)-regular sequence inside
  \(\mathfrak m\) whose length equals \(\text{spanFinrank}(\mathfrak m) = \dim(R)\). The proof is a
  strong induction on the dimension: the base case is the empty sequence, and the
  inductive step is \ref{mk-lem:regularLocal_inductive_step}. This is the lower-bound half
  of \ref{mk-cor:regular_cohen_macaulay}.
\end{lemma}
\begin{proof}
  \leanok

Use strong induction on \(\dim(R)\). The empty sequence proves the
  zero-dimensional case. In positive dimension,
  \ref{mk-lem:regularLocal_inductive_step} extends the regular sequence from
  the quotient of dimension one less, producing a sequence of length
  \(\dim(R)\).
\end{proof}

\section{Corollary: regular local rings are Cohen--Macaulay}
\label{mk-sec:ab_regular_cm}

\begin{definition}[Cohen--Macaulay local ring]\label{mk-def:cohen_macaulay_local}
  \lean{RingTheory.CohenMacaulay}
  \leanok
  \dcref{00N4,custom}

  \uses{mk-def:depth}
  \textit{Source: \cite{stacks-project}, tag 00N4 (definition-local-ring-CM).}
  A Noetherian local ring \((R, \mathfrak{m})\) is called \emph{Cohen--Macaulay} if
  \(\text{depth}(R) = \dim(R)\) (equivalently, if \(R\) is Cohen--Macaulay as a module
  over itself, where Cohen--Macaulay for a finite \(R\)-module \(M\) means
  \(\dim(\text{Supp}(M)) = \text{depth}(M)\), cf.\ Stacks tag 00N3).
\end{definition}

\begin{corollary}[Regular local rings are Cohen--Macaulay]\label{mk-cor:regular_cohen_macaulay}
  \lean{RingTheory.CohenMacaulay.of_regular}
  \leanok
  \dcref{00OD,custom}

  \uses{mk-def:cohen_macaulay_local, mk-thm:auslander_buchsbaum,
    mk-lem:exists_isRegular_of_regularLocal, mk-lem:length_le_ringKrullDim_of_isRegular}
  \textit{Source: \cite{stacks-project}, tag 00OD (lemma-regular-ring-CM).}
  Every regular Noetherian local ring is Cohen--Macaulay: if \((R, \mathfrak{m})\) is
  regular local of Krull dimension \(d\), then \(\text{depth}(R) = d = \dim(R)\).
\end{corollary}

\begin{proof}
  \leanok


Let \((R, \mathfrak{m})\) be regular of dimension \(d\) and pick a minimal generating set
  \(x_1, \dots, x_d\) of \(\mathfrak{m}\). By Stacks tag 00NQ (lemma-regular-domain), \(R\) is
  an integral domain, so \(x_1\) is a non-zero-divisor. The quotient \(R/x_1 R\) is a Noetherian
  local ring of dimension \(\geq d - 1\) (by Krull's principal ideal theorem,
  \texttt{lemma-one-equation}) with maximal ideal generated by \(x_2, \dots, x_d\), hence
  has dimension exactly \(d - 1\) and is again regular of dimension \(d - 1\). Inducting on \(d\),
  \(x_1, \dots, x_d\) is an \(R\)-regular sequence, so
  \(\text{depth}(R) \geq d\). The reverse inequality \(\text{depth}(R) \leq \dim(R) = d\)
  always holds for any nonzero finite module over a Noetherian local ring
  (Stacks tag 00LK, \texttt{lemma-bound-depth}). Thus \(\text{depth}(R) = d = \dim(R)\), i.e.\
  \(R\) is Cohen--Macaulay.

  \emph{Remark on the Auslander--Buchsbaum route.} The same conclusion can be obtained from
  \ref{mk-thm:auslander_buchsbaum} as follows: take \(M := R\) itself. Then \(M\) is
  finitely generated and projective (free of rank \(1\)), so \(\text{pd}_R(R) = 0\). The
  formula gives \(\text{depth}(R) = 0 + \text{depth}(R)\), a tautology. The non-tautological
  use of Auslander--Buchsbaum is rather: for a regular local ring of dimension \(d\), every
  finitely generated module \(M\) has \(\text{pd}_R(M) \leq d\) (Stacks tag 00O2,
  \texttt{proposition-finite-gl-dim-regular}); applying the formula to such an \(M\) then
  relates \(\text{depth}(M)\) to \(d - \text{pd}_R(M)\), which is the form most often used in
  practice. The direct regular-sequence argument above is the cleanest proof of
  \(R\) Cohen--Macaulay itself.
\end{proof}

\section{Application: Hartogs extension on a regular surface}
\label{mk-sec:ab_application_to_a4a}

Let \(S\) be a regular Noetherian surface and let \(x \in S\) be a closed point.  The local
Hartogs mechanism for the structure sheaf is:

\begin{enumerate}
  \item The local ring \(\mathcal{O}_{S,x}\) is a regular Noetherian local ring of Krull
    dimension \(2\) (because \(S\) is regular and \(\dim S = 2\)).
  \item By \ref{mk-cor:regular_cohen_macaulay} it is Cohen--Macaulay, so
    \(\text{depth}(\mathcal{O}_{S,x}) = 2\). Equivalently, the maximal ideal contains an
    \(\mathcal{O}_{S,x}\)-regular sequence of length \(2\).
  \item Sections of the structure sheaf \(\mathcal{O}_S\) on \(U = S \setminus \{x\}\) extend
    uniquely across \(x\) because \(H^0_x(\mathcal{O}_S) = 0\) and
    \(H^1_x(\mathcal{O}_S) = 0\) (these local-cohomology vanishings are equivalent to depth
    \(\geq 2\) at \(x\) by Stacks tag 0AVF; see Hartshorne, III.3.5, for the surface case).
\end{enumerate}

Thus \ref{mk-cor:regular_cohen_macaulay} supplies the depth-two input needed to extend regular
functions across a codimension-two point on a regular surface.  It does not extend a rational map
by choosing affine coordinates on the target: such an implication is false for general complete
targets.  The Albanese development instead combines the codimension-one valuative result with the
group-scheme theorem on indeterminacy loci; that geometric argument is independent of the
commutative-algebra result proved here.
